% !TeX program = lualatex
% !TeX root = main.tex

\documentclass[11pt,a4paper]{article}

\usepackage{fontspec}
\usepackage{polyglossia}
\setdefaultlanguage{english}
\setmainfont{CMU Serif}

\usepackage{amsmath}
\usepackage{amssymb}
\usepackage{booktabs}
\usepackage{float}
\usepackage{geometry}
\usepackage[hidelinks]{hyperref}
\usepackage[
    backend=biber,
    style=numeric,
    sorting=none
]{biblatex}

\addbibresource{references.bib}
\geometry{margin=25mm}

\title{Numerical Robustness of Restarted GMRES for Complex-Valued Linear Systems:\\
Preliminary Computational Experiments}
\author{Kirill Kabaev\\WaveCore Research}
\date{November 2026}

\begin{document}

\maketitle

\begin{abstract}
This preliminary study documents reproducible numerical checks of a sequential,
unpreconditioned restarted GMRES implementation for complex-valued sparse linear
systems. The implementation uses the Arnoldi process with a Hermitian inner
product, two passes of modified Gram--Schmidt orthogonalization, complex Givens
rotations, and an explicit true-residual check after every iteration. Small
deterministic systems probe nonnormality, ill-conditioning, restart-induced
stagnation, loss of orthogonality, extreme finite scaling, and disagreement
between the least-squares residual estimate and the recomputed residual. The
experiments show why residual accuracy and forward solution accuracy must be
reported separately. They also demonstrate, for the tested cases, that an
optimistic internal estimate does not produce a false convergence status. These
observations are verification evidence for one implementation; they are not a
claim of general numerical stability or scientific novelty.
\end{abstract}

\section{Introduction}

The generalized minimal residual method (GMRES) constructs approximations in a
Krylov subspace and minimizes a residual norm through a small least-squares
problem \cite{saad_schultz_1986}. Its use for nonsymmetric and complex-valued
systems makes implementation details---orthogonalization, least-squares updates,
stopping rules, and restarting---numerically consequential. The Arnoldi process
underlying GMRES originates in the construction of orthogonal Krylov bases
\cite{arnoldi_1951}; standard algorithmic treatments are given by
Saad \cite{saad_2003}.

This article records the first numerical observations from the WaveCore Research
GMRES implementation. The scope is deliberately narrow: deterministic dense test
matrices are converted to the project's complex CSR format and solved by the same
public routine intended for later sparse applications. Internal instrumentation
observes the actual Arnoldi basis without changing the returned solution or
status. The objectives are to verify implementation-level invariants, identify
failure modes before coupling the solver to a Helmholtz discretization, and retain
a traceable account of measured values.

No new GMRES variant or general stability theorem is proposed. In particular,
successful behavior on these small systems does not establish robustness for all
matrix classes, dimensions, restart lengths, or floating-point environments.

\section{Implemented restarted GMRES algorithm}

We solve
\begin{equation}
    A x = b, \qquad A\in\mathbb{C}^{n\times n},
    \quad x,b\in\mathbb{C}^{n},
\end{equation}
without preconditioning. Superscript $H$ denotes conjugate transpose, a bar denotes
scalar complex conjugation, and $\mathrm{i}^2=-1$. Vector norms are Euclidean:
$\lVert z\rVert_2=(\sum_{\ell=1}^n |z_\ell|^2)^{1/2}$. The implementation
accumulates norms with successive \texttt{hypot} calls rather than squaring all
entries directly. At the start of each restart cycle, let $x_{\mathrm{base}}$
denote the current approximation and set $r=b-Ax_{\mathrm{base}}$ and
$\beta=\lVert r\rVert_2$; in the first cycle $x_{\mathrm{base}}=x_0$ is the
user's initial guess. The stopping test is applied before normalization;
only if another iteration is required is $v_1=r/\beta$ formed. Complex vectors
use the Hermitian inner product
\begin{equation}
    \langle v,w\rangle = v^{H}w
    = \sum_{i=1}^{n}\overline{v_i}\,w_i.
    \label{eq:hermitian-product}
\end{equation}
The conjugation of the first argument in Eq.~\eqref{eq:hermitian-product} is
explicit in both the production implementation and the independent test
diagnostics.

Within a cycle, $j$ is the one-based Arnoldi step and
$V_j=[v_1,\ldots,v_j]\in\mathbb{C}^{n\times j}$. At step $j$, the implementation
initializes $w=Av_j$ and $h_{ij}=0$ for $1\leq i\leq j$. For each of two complete
modified Gram--Schmidt (MGS) passes, and sequentially for $i=1,\ldots,j$, it uses
\begin{align}
    \alpha &\leftarrow \langle v_i,w\rangle,\\
    h_{ij} &\leftarrow h_{ij}+\alpha,\qquad
    w\leftarrow w-\alpha v_i.
\end{align}
The same saved projection $\alpha$ is used in both updates. Set
$h_{j+1,j}=\lVert w\rVert_2$. If no near-breakdown is detected, the new vector is
$v_{j+1}=w/h_{j+1,j}$. In exact arithmetic, away from breakdown,
$AV_j=V_{j+1}\underline{H}_j$, where
$\underline{H}_j\in\mathbb{C}^{(j+1)\times j}$ is upper Hessenberg. The underline
denotes the rectangular Hessenberg matrix, not conjugation. This relation is
subject to rounding errors in the implementation. Reorthogonalization is a safeguard in
this implementation; finite-precision properties of MGS-GMRES and loss of
orthogonality are treated in detail by Paige, Rozlo\v{z}n\'{i}k, and
Strako\v{s} \cite{paige_rozloznik_strakos_2006}.

The Arnoldi relation gives the reduced least-squares problem
\begin{equation}
    y_j = \arg\min_{y\in\mathbb{C}^j}
    \left\lVert \beta e_1-\underline{H}_j y\right\rVert_2,
    \qquad x_j=x_{\mathrm{base}}+V_j y_j.
    \label{eq:least-squares}
\end{equation}
Here $e_1\in\mathbb{R}^{j+1}$ is the first coordinate vector, and $x_j$ is the
candidate within the current cycle. The upper-Hessenberg matrix is reduced by
left multiplication with complex unitary Givens rotations
\begin{equation}
    G = \begin{bmatrix}c&s\\-s^{\ast}&c\end{bmatrix},
    \qquad c\in\mathbb{R},\quad s\in\mathbb{C},\quad c^2+|s|^2=1.
    \label{eq:givens}
\end{equation}
The scalar $s^{\ast}=\overline{s}$ is explicitly conjugated; thus $G^H G=I_2$.
For entries $a,b\in\mathbb{C}$ to be rotated, the code sets
$d=\operatorname{hypot}(|a|,|b|)$ and, when both entries are nonzero,
\begin{equation}
    c=\frac{|a|}{d},\qquad
    s=\frac{a}{|a|}\frac{\overline{b}}{d},\qquad
    G\begin{bmatrix}a\\b\end{bmatrix}
      =\begin{bmatrix}(a/|a|)d\\0\end{bmatrix}.
\end{equation}
If $b=0$, it uses $(c,s)=(1,0)$, including $a=b=0$. If $a=0$ and
$b\ne0$, the code returns
\[
    c=0,\qquad s=\frac{\overline{b}}{|b|},\qquad
    G\begin{bmatrix}0\\b\end{bmatrix}
      =\begin{bmatrix}sb\\0\end{bmatrix}
      =\begin{bmatrix}\overline{b}b/|b|\\0\end{bmatrix}
      =\begin{bmatrix}|b|\\0\end{bmatrix}.
\]
The conjugate in $s$ fixes the phase of the first entry to the nonnegative real
number $|b|$. For example, $b=\mathrm{i}$ gives $s=-\mathrm{i}$ and
$s b=(-\mathrm{i})\mathrm{i}=1$. In contrast, $s=b/|b|$ would produce
$b^2/|b|$ (equal to $-1$ for $b=\mathrm{i}$) and is not the convention used by
\texttt{make\_rotation} and \texttt{rotate}.
Previous rotations are applied to each new Hessenberg column, and the new
rotation is applied both to that column and to the reduced right-hand side.
Triangular back substitution then solves the nonsingular reduced system.

A restart cycle allows at most $\min(m,n,N_{\max}-N)$ Arnoldi steps, where $m$
is the requested restart length, $N$ is the total completed iteration count at
cycle entry, and $N_{\max}$ is the global limit. Convergence or breakdown may
end a cycle earlier. The last approximation of a completed cycle becomes the
next cycle's base point.

\subsection{Stopping and diagnostic rules}

Use a common normalization for recomputed and estimated residuals:
\begin{equation}
    D_b=\begin{cases}\lVert b\rVert_2,&b\ne0,\\1,&b=0,\end{cases}
    \qquad \rho(x)=\frac{\lVert b-Ax\rVert_2}{D_b}.
    \label{eq:true-residual}
\end{equation}
For $b\ne0$, this is relative; for the exact zero right-hand side it is absolute.
Let $g^{(j)}$ be the reduced vector $\beta e_1$ after the first $j$ embedded
Givens rotations. The least-squares residual estimate is
\begin{equation}
    \widehat{\rho}_j=
    \frac{|g^{(j)}_{j+1}|}{D_b}.
\end{equation}
The solver never uses this estimate to declare convergence. It recomputes
$b-Ax_j$ in double precision with a CSR product after each candidate update and
returns \texttt{CONVERGED} only if the computed $\rho(x_j)\leq\tau$, where
$\tau>0$ is the finite user tolerance. The same check at the initial guess can
return success with zero iterations. These computed residuals, including the
initial value, form the public residual history; they are not exact-arithmetic
residuals or the independent wider-precision test values.

Near-breakdown is flagged when
\begin{equation}
    \lVert w\rVert_2 \leq
    16\,\epsilon_{\mathrm{double}}\,\lVert Av_j\rVert_2.
\end{equation}
Here $w$ is the remainder after both MGS passes, and
$\epsilon_{\mathrm{double}}=2^{-52}$ is the spacing above one in the reported
binary64 environment. If back substitution succeeds, convergence is checked
before this flag terminates the solve: a residual meeting the tolerance gives
happy breakdown; otherwise the status is \texttt{BREAKDOWN}. An exactly zero
diagonal in the rotated triangular problem also returns \texttt{BREAKDOWN},
retaining the previous approximation and repeating its residual in the history.
Exhausting $N_{\max}$ returns \texttt{NOT\_CONVERGED}. Invalid dimensions or
options are rejected, and detected nonfinite arithmetic raises an exception
rather than returning success.

\section{Computational methodology}

The reported values were reproduced from the deterministic CTest executable
\nolinkurl{wavecore_gmres_test} at Git revision
\nolinkurl{26bf63a73a7201c36d6e6d6933942624ef74d9de}. The run used GCC~13.3.0 and
CMake~3.28.3 on x86-64 Ubuntu under WSL2. The production vectors store
double-precision C++ complex numbers (53 binary significand bits). The build uses
C++20 in Debug mode, with \texttt{-g -Wall -Wextra -Wpedantic} and no explicit
optimization flag. Test-side residuals, solution errors, and basis inner products
for the extended diagnostic systems are accumulated independently with
\texttt{std::complex<long double>} before conversion to the reported decimal
values; \texttt{long double} has 64 significand bits on this platform. Baseline
tests instead recompute residuals with independent dense double arithmetic.
All computations are sequential and deterministic; no random seed is required.
The revision identifies the solver, diagnostics, tests, and CMake sources, not a
committed version of this manuscript.

For a prescribed reference vector $x_\star$, each right-hand side is formed by
the independent dense routine in double precision:
$b=\operatorname{fl}(Ax_\star)$. Thus $x_\star$ is the intended solution of the
mathematical construction; it need not be the exact solution of the stored
rounded system. The independently evaluated relative reference error is
\begin{equation}
    \eta_x=\frac{\lVert x-x_\star\rVert_2}{\lVert x_\star\rVert_2},
\end{equation}
and the independently evaluated residual uses Eq.~\eqref{eq:true-residual} with
the stored $A$, $b$, and computed $x$. The wide oracle reduces arithmetic error
in verification; it does not recover unrounded inputs. In the tables, $\rho$
denotes this independent residual, whereas $\rho_{\mathrm{solver}}$ denotes the
double-precision value used by the solver's stopping test.

At observer sample $t$, after a nonsingular reduced solve, define the diagnostic
matrix $W_t$ as the current $V_j$ augmented by $w/\lVert w\rVert_2$ if no
near-breakdown is detected; otherwise set $W_t=V_j$. This includes the prospective
next vector even when the iteration has just converged or ended a cycle. With
$q_t$ columns in $W_t$, the diagnostic records
\begin{equation}
    \delta_V=\max_t\lVert W_t^H W_t-I_{q_t}\rVert_F,
    \qquad \lVert M\rVert_F=
        \left(\sum_{a,b}|M_{ab}|^2\right)^{1/2}.
    \label{eq:orthogonality}
\end{equation}
Instrumentation is internal to the tests. A paired assertion verifies that it
does not change the public solver's status, iterations, solution, or residual
history. The maximum spans all observer samples and restart cycles, not the
public residual-history index. The one-pass control instead reports the final
defect of its nine-vector basis after eight Arnoldi steps; it does not take a
maximum over iterations. The reported maximum residual-estimate gap is
$\max_t|\widehat{\rho}_t-\rho_{\mathrm{solver},t}|$, using the double-precision
recomputed residual at each observer call, not the final wide oracle.

\subsection{Matrices and solver parameters}

The principal diagnostic systems are as follows.
\begin{itemize}
    \item \textbf{Complex Hilbert(10):}
    $A_{ij}=\exp(0.3\mathrm{i}(i-j))/(i+j+1)$ for zero-based $i,j$, a unitary
    diagonal similarity of the real Hilbert matrix in exact arithmetic. The reference is
    $(x_\star)_i=1+0.25\mathrm{i}(i\bmod 3)$.
    \item \textbf{Near-singular diagonal:}
    $A=\operatorname{diag}(1,10^{-14})$ and $x_\star=(1,1)^T$. The stated
    mathematical matrix has exact 2-norm condition number $10^{14}$.
    \item \textbf{Complex nonnormal Jordan-like matrix:} a $6\times6$ upper
    bidiagonal matrix with diagonal $1+0.5\mathrm{i}$ and superdiagonal
    $2-\mathrm{i}$; $(x_\star)_i=(i+1)-0.25\mathrm{i}$.
    \item \textbf{Scaled complex matrix:} $A=\sigma B$ for
    $\sigma\in\{10^{-200},1,10^{200}\}$, where
    \begin{equation}
        B=\begin{bmatrix}
            4+\mathrm{i}&1-2\mathrm{i}&1/2\\
            \mathrm{i}&5-\mathrm{i}&-1+\mathrm{i}/2\\
            2&-\mathrm{i}&6+2\mathrm{i}
        \end{bmatrix},\qquad
        x_\star=\begin{bmatrix}1+2\mathrm{i}\\-2+\mathrm{i}/2\\1/4-\mathrm{i}\end{bmatrix}.
    \end{equation}
    The residual-gap experiment uses $B$ and this same reference vector.
    \item \textbf{Restart stagnation:}
    $A=\left[\begin{smallmatrix}0&-1\\1&0\end{smallmatrix}\right]$,
    $b=(1,0)^T$, and $x_\star=(0,-1)^T$. GMRES(1) stagnates exactly at the zero
    approximation for this initial residual.
\end{itemize}

The restart-one baseline uses
\begin{equation}
    T=\begin{bmatrix}
        2+\mathrm{i}&0.2-0.1\mathrm{i}&0\\
        0&3-0.5\mathrm{i}&0.1+0.2\mathrm{i}\\
        0&0&4+\mathrm{i}
    \end{bmatrix},\qquad
    x_\star=\begin{bmatrix}1-\mathrm{i}\\2+0.5\mathrm{i}\\-1+2\mathrm{i}\end{bmatrix}.
\end{equation}
The restart-two diagonal baseline uses $\operatorname{diag}(2,3,4,5)$ and
$x_\star=(1+\mathrm{i},2-\mathrm{i},-1+0.5\mathrm{i},0.5+2\mathrm{i})^T$.
Both baseline tests initialize every component to $0.125-0.25\mathrm{i}$.
The inconsistent zero-operator test uses a $2\times2$ zero matrix,
$b=(1,1)^T$, and a zero initial guess. The following options reproduce all
tabulated systems and the residual-gap experiment; the default diagnostic
tolerance is $10^{-12}$ except for the indicated gap test.

\begin{table}[H]
\centering
\small
\caption{Options in the executable tests; $m$ is the requested restart length.}
\label{tab:options}
\begin{tabular}{lrrrr}
\toprule
System & $n$ & $m$ & $N_{\max}$ & $\tau$ \\
\midrule
Complex Hilbert(10) & 10 & 10 & 30 & $10^{-12}$ \\
Near-singular diagonal & 2 & 2 & 4 & $10^{-12}$ \\
Nonnormal Jordan-like & 6 & 6 & 12 & $10^{-12}$ \\
Scaled $B$ (each scale) & 3 & 99 & 6 & $10^{-12}$ \\
$B$, residual-gap test & 3 & 3 & 6 & $10^{-30}$ \\
$T$, restart-one baseline & 3 & 1 & 100 & $10^{-12}$ \\
Diagonal restart-two baseline & 4 & 2 & 100 & $10^{-12}$ \\
Skew system & 2 & 1 & 5 & $10^{-12}$ \\
Skew system, limit checks & 2 & 2 & 1 or 2 & $10^{-12}$ \\
Inconsistent zero operator & 2 & 2 & 10 & $10^{-12}$ \\
\bottomrule
\end{tabular}
\end{table}

All cases in Table~\ref{tab:options} start at zero except the two explicitly
identified baseline tests. The one-pass Hilbert control starts from the normalized
same right-hand side and performs eight Arnoldi steps without QR updates or
GMRES stopping tests. It uses a dense product and a norm accumulated as
$\operatorname{hypot}(\cdot,|z_i|)$, while production uses CSR and separate real
and imaginary components in its norm. Thus it is a diagnostic comparison, not
an isolated experiment changing only the MGS pass count.

Other tests cover identity and permutation matrices, an empty system, a zero
right-hand side, exact initial guesses, invalid arguments, and finite inputs
whose matrix--vector product overflows.

\subsection{Reproduction commands}

From the repository root containing the stated solver/test revision, run:
\begin{verbatim}
git rev-parse HEAD
git diff 26bf63a73a7201c36d6e6d6933942624ef74d9de -- \
  CMakeLists.txt include/wavecore/helmholtz/gmres.hpp \
  src/helmholtz/gmres.cpp src/helmholtz/gmres_diagnostics.hpp \
  tests/helmholtz/test_gmres.cpp
cmake -S . -B build -DCMAKE_BUILD_TYPE=Debug \
  -DBUILD_TESTING=ON
cmake --build build --parallel 2
ctest --test-dir build --output-on-failure
set -o pipefail
./build/wavecore_gmres_test | tee /tmp/wavecore-gmres-results.txt
\end{verbatim}
The first two commands record provenance and reveal changes to the numerical
sources. Record compiler, CMake, and platform versions alongside the saved
stdout. The diagnostic output labels \texttt{error}, \texttt{wide\_residual},
\texttt{estimated}, \texttt{orthogonality\_F}, and \texttt{max\_residual\_gap}
map to the quantities defined above. Baseline lines use \texttt{true\_residual}
for the independent dense double residual. Printed status codes 0, 1, and 2 mean
\texttt{CONVERGED}, \texttt{NOT\_CONVERGED}, and \texttt{BREAKDOWN}, respectively.
Small changes in the last digits across platforms remain possible.

The manuscript requires the existing LuaLaTeX/Biber installation:
\begin{verbatim}
cd papers/02_gmres_numerical_observations
latexmk -lualatex -interaction=nonstopmode \
  -file-line-error -halt-on-error main.tex
\end{verbatim}
The root \texttt{make ci} currently verifies the first paper and MkDocs website;
it does not replace the separate build of this manuscript or CTest.

\section{Measured results}

Table~\ref{tab:principal} reports independently recomputed errors and residuals.
Status abbreviations are C for \texttt{CONVERGED}, NC for
\texttt{NOT\_CONVERGED}, and B for \texttt{BREAKDOWN}. All displayed digits are
rounded from the deterministic test output; they should not be interpreted as
uncertainty estimates.

\begin{table}[H]
\centering
\small
\caption{Principal diagnostic systems.}
\label{tab:principal}
\begin{tabular}{lrrrrr}
\toprule
System & Iter. & Status & $\eta_x$ & $\rho$ & $\delta_V$ \\
\midrule
Complex Hilbert(10) & 9 & C
    & $5.473\times10^{-2}$ & $5.610\times10^{-15}$ & $5.453\times10^{-16}$ \\
Near-singular diagonal & 1 & C
    & $7.071\times10^{-1}$ & $1.000\times10^{-14}$ & $0$ \\
Nonnormal Jordan-like & 6 & C
    & $2.860\times10^{-15}$ & $1.368\times10^{-16}$ & $4.526\times10^{-16}$ \\
GMRES(1) stagnation & 5 & NC
    & $1$ & $1$ & $0$ \\
\bottomrule
\end{tabular}
\end{table}

For Complex Hilbert(10), a local one-pass MGS control built only for the test
produced a final nine-vector defect of $1.022\times10^{-4}$ after eight Arnoldi
steps. The actual two-pass production basis reached a maximum defect of
$5.453\times10^{-16}$ in the corresponding diagnostic solve. This comparison is
specific to this matrix and execution. The one-pass endpoint and two-pass maximum
have different aggregation rules, and the controls also differ in product and
norm implementation as described above. They do not prove a universal bound for
two-pass MGS.

\begin{table}[H]
\centering
\small
\caption{Finite scaling experiment; restart was requested as 99 and capped at
$n=3$.}
\label{tab:scaling}
\begin{tabular}{rrrrrr}
\toprule
Scale & Iter. & $\eta_x$ & $\rho$ & $\widehat{\rho}$ & $\delta_V$ \\
\midrule
$10^{-200}$ & 3 & $4.276\times10^{-16}$ & $2.716\times10^{-16}$
    & $1.534\times10^{-33}$ & $2.076\times10^{-16}$ \\
$1$ & 3 & $1.711\times10^{-16}$ & $1.608\times10^{-16}$
    & $4.284\times10^{-32}$ & $6.618\times10^{-16}$ \\
$10^{200}$ & 3 & $1.763\times10^{-16}$ & $1.443\times10^{-16}$
    & $2.995\times10^{-33}$ & $2.913\times10^{-16}$ \\
\bottomrule
\end{tabular}
\end{table}

The three scaling cases in Table~\ref{tab:scaling} converged in three iterations
with the requested tolerance $10^{-12}$. These observations establish coverage of
the selected finite scales, not scale invariance for arbitrary matrices or data.

\begin{table}[H]
\centering
\small
\caption{Stopping and restart boundary cases.}
\label{tab:stopping}
\begin{tabular}{lrrrr}
\toprule
Case & Restart & $N_{\max}$ & Iter. & Status / final $\rho$ \\
\midrule
Complex $3\times3$ baseline & 1 & 100 & 27 & C / $3.803\times10^{-13}$ \\
Diagonal $4\times4$ baseline & 2 & 100 & 19 & C / $7.170\times10^{-13}$ \\
Skew system & 1 & 5 & 5 & NC / $1$ \\
Skew system & 2 & 1 & 1 & NC / $1$ \\
Skew system & 2 & 2 & 2 & C / $0$ \\
Inconsistent zero operator & 2 & 10 & 1 & B / $1$ \\
\bottomrule
\end{tabular}
\end{table}

The first two rows and the zero-operator row report independently computed dense
double residuals; the skew rows use the wide diagnostic oracle. A repeated run
reproduced all displayed measured values; the dimension and iteration-limit
metadata above follow the executable test options.
Table~\ref{tab:stopping} confirms that the total count, rather than a per-restart
count, is compared with $N_{\max}$. The exact GMRES(1) stagnation case retains
$\rho=1$ in every history entry and terminates without a false success.

\section{Discussion}

\subsection{Residual accuracy is not solution accuracy}

The most important observation is the separation between a small residual and a
small forward error. For the near-singular diagonal system, the component
associated with $10^{-14}$ contributes only $10^{-14}$ to the unnormalized
residual when it is unresolved. With a tolerance of $10^{-12}$, the solver
therefore returns convergence after one iteration even though
$\eta_x=0.7071$. This is expected conditioning behavior, not a contradiction of
the stopping rule.

The complex Hilbert experiment is less extreme but qualitatively similar:
$\rho=5.610\times10^{-15}$ accompanies a $5.47\%$ relative solution error. A
residual is a backward-error-like observable; it is not the usual normwise
backward error $\lVert b-Ax\rVert_2/(\lVert A\rVert_2\lVert x\rVert_2+
\lVert b\rVert_2)$. For nonsingular $A$, nonzero $x_\star$, and $b=Ax_\star$,
\begin{equation}
    \eta_x\leq\kappa_2(A)\rho(x),\qquad
    \kappa_2(A)=\lVert A\rVert_2\lVert A^{-1}\rVert_2.
\end{equation}
For the tests' rounded right-hand side, write $\delta b=b-Ax_\star$; then
\begin{equation}
    \lVert x-x_\star\rVert_2\leq\lVert A^{-1}\rVert_2
       \bigl(\lVert b-Ax\rVert_2+\lVert\delta b\rVert_2\bigr).
\end{equation}
Both the residual and right-hand-side formation error can therefore be amplified
by conditioning. General finite-precision and conditioning principles are
reviewed by Higham \cite{higham_2002}. Consequently,
future experiment reports must retain both $\rho$ and $\eta_x$ whenever a trusted
reference solution is available.

\subsection{Estimated and true residuals}

The small least-squares problem can yield a much more optimistic residual estimate
than a fresh evaluation of $b-Ax$. For the nonnormal Jordan-like system, the last
estimate was $6.079\times10^{-33}$ while the independently accumulated true
residual was $1.368\times10^{-16}$. In the unscaled $3\times3$ system, the
corresponding pair was $4.284\times10^{-32}$ and
$1.608\times10^{-16}$.

The latter system was repeated with a deliberately stringent tolerance
$10^{-30}$. The estimate crossed the tolerance, but the true residual did not;
the solver returned \texttt{BREAKDOWN} after three iterations rather than
\texttt{CONVERGED}. This tolerance was not attained in the reported run; it is
not claimed to be unattainable for every double-precision system. The maximum
absolute gap between the estimate and the solver's double-precision residual,
$\max_t|\widehat{\rho}_t-\rho_{\mathrm{solver},t}|$, was
$1.665\times10^{-16}$. The independent wide residual is the final value quoted
above, not a history recorded at each step. Thus, in this reproducible example,
the independent residual guard prevents false convergence. It does not eliminate
roundoff from the true-residual computation itself, which remains a double-
precision CSR matrix--vector product.

\subsection{Restart and breakdown behavior}

Restart length changes the available minimization space. The skew-symmetric
$2\times2$ example demonstrates exact GMRES(1) stagnation, whereas restart two
recovers the exact solution in two iterations. Other baseline systems converge
with restart one or two, but require 27 and 19 total iterations, respectively.
These cases verify cycle continuation and global iteration accounting without
supporting a general convergence-rate claim.

The implementation distinguishes a true-residual success at near-breakdown from
breakdown without convergence. Identity-like cases exercise happy breakdown; the
inconsistent zero-operator case returns \texttt{BREAKDOWN} with residual one.
This status communicates that the Krylov process cannot supply another useful
direction under the implemented threshold.

\section{Limitations and planned Helmholtz integration}

The present evidence has the following explicit limitations.
\begin{itemize}
    \item The matrices are small, deterministic diagnostic examples. They do not
    represent the dimensions, spectra, nonnormality, or memory traffic of a
    production sparse wave problem.
    \item Only one compiler/platform run is reported. No cross-platform rounding
    study, sanitizer run, performance benchmark, or statistical ensemble is
    included.
    \item The solver is sequential and unpreconditioned. No claim is made about
    practical iteration counts for indefinite Helmholtz systems.
    \item Orthogonality is observed by internal test instrumentation and checked
    with extended-precision accumulation, but the experiment is not a formal
    rounding-error analysis.
    \item The true residual is recomputed in double precision. Severe cancellation,
    underflow, or overflow can still limit its reliability; detected nonfinite
    arithmetic is rejected, but all finite extreme cases cannot be covered.
    \item A small residual alone cannot certify a small solution error for an
    ill-conditioned system. Condition estimation is not implemented.
\end{itemize}

\subsection{Data required for iteration plots}

The current tests print the final residual estimate, the final independent wide
residual, the maximum estimate--solver-residual gap, and the maximum two-pass
orthogonality defect. The public result retains the solver's residual history
only in memory. The one-pass control returns a single final defect. These
aggregates do not determine iteration curves, so no such curves are inferred
from them in this manuscript.

An estimate-versus-true-residual plot would require exporting each observer
sample as a row containing the case, total iteration, restart cycle, local step,
$\widehat{\rho}_t$, $\rho_{\mathrm{solver},t}$, basis size, and breakdown flag.
A separate per-step wide residual would require access to each candidate
solution, which the present observer does not expose. The plot must label the
double-precision recomputation and the wide oracle distinctly.

An orthogonality plot would require collecting
$\lVert W_t^H W_t-I\rVert_F$ at every step of the production run and of the
local one-pass control, using the same Hilbert matrix, right-hand side, step
range, and basis-size convention. To isolate the effect of the MGS pass count,
the product and norm routines should also be matched. Raw CSV data and the
revision, compiler, flags, precision, and solver options should accompany any
figures. Zero values must be shown explicitly or identified as omitted on
logarithmic axes, not replaced by invented positive measurements. Producing
these traces requires additional test instrumentation and a repeated run; it
does not require changing the numerical solver algorithm.

\subsection{Integration with FDFD Helmholtz}

The next application target is the WaveCore five-point FDFD discretization of
$-\Delta u-k^2u=f$ with homogeneous Dirichlet boundaries. Integration should
first use manufactured solutions away from both continuous and discrete
resonances. Each run should record grid dimensions, $k$, restart length, tolerance,
iteration limit, solver status, complete residual history, independently
recomputed algebraic residual, and error against the manufactured solution.
Grid-refinement studies must separate discretization error from GMRES iteration
error. Near resonances, the experiment should report conditioning indicators and
must not infer field accuracy from residual size alone.

Only after this unpreconditioned baseline is verified should preconditioning be
introduced. Candidate preconditioners and performance changes will require
separate hypotheses, numerical-equivalence checks, and measurements; none are
evaluated in the present article.

\section{Conclusions}

The tested WaveCore implementation is consistent with its stated algorithm on the
reported deterministic systems: complex inner products are Hermitian, Arnoldi
uses two MGS passes, restart iterations obey one global limit, and convergence is
accepted only after recomputing the true residual. The diagnostics expose three
practical distinctions that must remain visible in later work: residual versus
forward error, least-squares estimate versus recomputed residual, and stagnation
versus convergence.

Two results are particularly relevant for future wave simulations. First, a
relative residual near machine precision can coexist with a percent-level
solution error on an ill-conditioned problem. Second, an internal estimate far
below machine precision did not cause a false success because the implementation
checked $b-Ax$ explicitly. These are preliminary observations from selected
tests, not proofs of robustness. FDFD Helmholtz integration, conditioning studies,
larger sparse experiments, and preconditioner evaluation remain future work.


\printbibliography

\end{document}
